% -----------------------------------------------------------------------------
% common header for all exercise sheets
\documentclass[
paper=a4,
parskip=half,
oneside
]{scrartcl}


\subject{Hanoi Fall School 2013 on Numerical Analysis}
\author{}
\date{}

\usepackage[USenglish]{babel}

\usepackage{graphicx}
% The official corporate design guidelines
% <http://www.pressestelle.tu-berlin.de/online_styleguide/design/typografie_ikonografie/>
% advise to set the logo flush right and at 22% of the pagewidth. Be somewhat
% more conservative with the size here.
%
% Use
\usepackage{pbox}
% to getting an authors box that is exactly as wide as the longest line in it,
% cf. <https://www.tug.org/tugboat/tb25-2/tb81law.pdf>.
\titlehead{
\pbox{\textwidth}{Prof.\ Dr.\ G\"unter B\"arwolff\\
                  Dr.\ Antoine Laurain\\
                  Dr.\ Christian Schr\"oder\\
                  Tobias Ahnert \\
                  Ute Kandler\\
                  Nico Schl\"omer}
\hfill
%\includegraphics[height=0.1\textwidth]{logos/bms-logo-eff}
\includegraphics[height=0.1\textwidth]{logos/bms-logo-bw}
\hfill
%\includegraphics[width=0.15\paperwidth]{logos/tu-logo-long}
\includegraphics[height=0.1\textwidth]{logos/tu_logo_short_improved}
}

% We're not providing \author or \date, so define a new \maketitle that
% ignores those (and doesn't reserve any extra space), cf.
% <http://tex.stackexchange.com/questions/134851/when-author-date-fields-blank-remove-whitespace-between-header-and-text>.
\makeatletter
\renewcommand{\@maketitle}{%
  \clearpage
  \let\footnote\thanks
  \ifx\@extratitle\@empty \else
    \noindent\@extratitle \next@tpage \if@twoside \null\next@tpage \fi
  \fi
  \setparsizes{\z@}{\z@}{\z@\@plus 1fil}\par@updaterelative
  \ifx\@titlehead\@empty \else
    \begin{minipage}[t]{\textwidth}
      \@titlehead
    \end{minipage}\par
  \fi
  \null
  \vskip 2em%
  \begin{center}%
    \ifx\@subject\@empty \else
      {\subject@font \@subject \par}
      \vskip 1.5em
    \fi
    {\titlefont\huge \@title \par}%
    \vskip .5em
    {\ifx\@subtitle\@empty\else\usekomafont{subtitle}\@subtitle\par\fi}%
    %\vskip 1em
    %{\Large
    %  \lineskip .5em%
    %  \begin{tabular}[t]{c}
    %    \@author
    %  \end{tabular}\par
    %}%
    %\vskip 1em%
    %{\Large \@date \par}%
    %\vskip \z@ \@plus 1em
    %{\Large \@publishers \par}
    %\ifx\@dedication\@empty \else
    %  \vskip 2em
    %  {\Large \@dedication \par}
    %\fi
  \end{center}%
  \par
  \vskip 2em
}%
\makeatother



% use the latin modern font
\usepackage[T1]{fontenc}
\usepackage{lmodern}

% use microtype
\usepackage{microtype}

\usepackage[
bookmarks,
colorlinks
]{hyperref}

% define MATLAB(R) symbol
\newcommand\matlab{{MATLAB\textsuperscript{\textregistered}}}

% redefine the enumerate environment in such a way that the pattern looks like
%
% (a)
% (b) (i)
%     (ii)
% (c)
%
\renewcommand\theenumi   {\alph{enumi}}
\renewcommand\labelenumi {(\theenumi)}
\renewcommand\theenumii  {\roman{enumii}}
\renewcommand\labelenumii{(\theenumii)}

% Some common definitions.
\usepackage{amsmath}
\usepackage{amsfonts}
\newcommand\N{\ensuremath{\mathbb{N}}}
\newcommand\R{\ensuremath{\mathbb{R}}}
\newcommand\dd{\ensuremath{\text{d}}}
\newcommand\tp{\ensuremath{\text{\upshape T}}}
\newcommand\ii{\ensuremath{\text{\upshape i}}}
\newcommand\dfn{\ensuremath{\mathrel{\mathop:}=}}


\usepackage{algpseudocode}

\title{Final Examination}

\newtheorem{algorithm}{Algorithm}
% -----------------------------------------------------------------------------
\begin{document}

\maketitle

\section*{Solving linear equations}

\paragraph{Problem 1}

\paragraph{Problem 2}
% -----------------------------------------------------------------------------
\paragraph{Problem 3}
% -----------------------------------------------------------------------------
\section*{Nonlinear optimization}
% -----------------------------------------------------------------------------
\paragraph{Problem 4 (Definitions)}
See syllabus.
% -----------------------------------------------------------------------------
\paragraph{Problem 5 (Conjugate gradient method)}
Let
\[
    f: \R^n \to \R,\quad f(x) = \frac{1}{2} x^\tp A x - b^\tp x,
\]
with $A\in\R^{n\times n}$, symmetric and positive definite, $b\in\R^n$.

\begin{enumerate}
  \item We have that
\[
\nabla f(x) = Ax - b,
\]
and since $A$ is invertible (because it is positive definite), the only
critical point of $f$ is $x^*=A^{-1}b$.

The Hessian of $f$, $\nabla^2 f(x) = A$, is obviously positive definite for
$x^*$, so the critical point is a minimum.

  \item See syllabus, page 36;
\[
\alpha_k = \frac{\|\nabla f(x_k)\|}{\|\nabla f(x_k)\|_A}.
\]

  \item See syllabus, page 41;
\[
\beta_i^k = \frac{g_{k+1}^\tp A d_i}{d_i^\tp A d_i} \quad\forall i\in\{1,\dots,k\}.
\]
  The directions $d_k$ are called \emph{conjugate} or \emph{$A$-conjugate}.
\end{enumerate}
% -----------------------------------------------------------------------------
\paragraph{Problem 6 (Newton's method)}

Let
\[
  f: \R \to \R,\quad f(x) \dfn 1 - \exp(-x^2).
\]

\begin{enumerate}

\item We have
\[
\nabla f(x) = 2x\exp(-x^2),
\]
so the only critical point is $x^*=0$.

Moreover, it clearly is $f(x)\ge 0$, and $f(0) = 0$, so it must be a global
minimizer.

\item For which $x$ is $\nabla^2 f(x)$ positive definite?
\[
\nabla^2 f(x) = 2\exp(-x^2) - 4x^2\exp(-x^2) = (2-4x^2)\exp(-x^2),
\]
so $\nabla^2 f(x)$ is positive for $\{x\in\R: |x|<\sqrt{1/2}\}$.

\item One Newton step is defined by
\[
x_{k+1} = x_k - [\nabla^2 f(x_k)]^{-1} \nabla f(x_k),
\]
so in this case
\begin{multline*}
x_{k+1}
= x_k - \frac{2x_k\exp(-x_k^2)}{(2-4x_k^2)\exp(-x_k^2)}
= x_k - \frac{x_k}{(1-2x_k^2)}\\
= \left[1 - \frac{1}{1-2x_k^2}\right] x_k
= \left[\frac{-2x_k^2}{1-2x_k^2}\right] x_k
= \underbrace{\left[\frac{1}{1-\frac{1}{2x_k^2}}\right]}_{=\gamma} x_k.
\end{multline*}
We ask ourselves the question when
\begin{equation}\label{eq:gamma}
1 > |\gamma| = \left|\frac{1}{1-\frac{1}{2x_k^2}}\right|
\end{equation}
since then, the iteration converges to the minimum $x^*=0$.
We separate two cases for $\gamma$.
\begin{itemize}
\item $1-\frac{1}{2x_k^2} > 0$, i.e., $x_k^2 > \frac{1}{2}$.
      In this case, inequality (\ref{eq:gamma}) is equivalent to
\[
1 < 1 - \frac{1}{2x_k^2},
\]
i.e.,
\[
0 < - \frac{1}{2x_k^2},
\]
a contradiction.

\item $1-\frac{1}{2x_k^2} < 0$, i.e., $x_k^2 < \frac{1}{2}$.
      Here, inequality (\ref{eq:gamma}) is equivalent to
\[
1 < -1 + \frac{1}{2x_k^2},
\]
i.e.,
\[
x_k^2 < \frac{1}{4}.
\]
Hence, as soon as $|x_k| < \frac{1}{2}$, the $|\gamma| < 1$ and the iteration
converges towards $x^*=0$.

We now turn to the quotient convergence factor $Q_p(\{x_k\})$. We have
\begin{multline*}
Q_p(\{x_k\})
= \limsup_{k\to\infty} \frac{\|x_{k+1}-x^*\|}{\|x_k-x^*\|^p}
= \limsup_{k\to\infty} \frac{\left|\frac{1}{1-1/(2x_k^2)} x_k\right|}{|x_k|^p}\\
= \limsup_{k\to\infty} \frac{1}{-1+1/(2x_k^2)} |x_k|^{1-p}
= \limsup_{k\to\infty} \frac{2}{1-2x_k^2} |x_k|^{3-p},
\end{multline*}
so $Q_p(\{x_k\})=0$ for any $p\in[1,3)$, $Q_3(\{x_k\}) = 2$, and otherwise nonexistent.
\end{itemize}



  \item Let
    \[
       g: \R^2 \to \R,\quad g(x_1, x_2) \dfn x_1^4 + x_2^2.
    \]
    Do tasks (a), (b), (c) for $g$.
 How do you explain the difference in the Q-convergence order between $f$ and $g$?
\end{enumerate}
% -----------------------------------------------------------------------------
\section*{Numerical methods for PDEs}
% -----------------------------------------------------------------------------
\paragraph{Problem 7}
% -----------------------------------------------------------------------------
\paragraph{Problem 8}
% -----------------------------------------------------------------------------
\paragraph{Problem 9*}
% -----------------------------------------------------------------------------


\end{document}
